%!TEX root = ../main.tex
% =============================================================
%  CAPÍTULO 3 — MARCO TEÓRICO Y ESTADO DEL ARTE
%                                       (Entrega 3 · 10 % · 28/7)
% =============================================================

\chapter{Síntesis de la literatura consultada}\label{cap:marco-teorico}

\begin{guia}[Guía de la cátedra: qué se evalúa en este capítulo]
\begin{enumerate}
    \item \textbf{Estado del arte}: investigaciones nacionales,
          internacionales y productos similares, por separado. De cada
          trabajo: qué hizo y qué resultados obtuvo, con métricas.
    \item \textbf{Brechas identificadas}: qué falta en lo existente y
          justifica \emph{este} proyecto. Es lo que separa un estado del
          arte de un listado. Incluir también los resultados negativos de
          la literatura.
    \item \textbf{Bases teóricas}: desarrollan las variables del título y su
          relación. Cada concepto termina explicando cómo se usa en el
          proyecto (teoría instrumental, no de relleno). Citas canónicas
          antes que tutoriales.
    \item \textbf{Marco conceptual}: conceptos que se nombran pero no se
          desarrollan. Los términos del glosario se marcan con
          \verb|\gls{}|, por ejemplo \gls{g:mvp} y \gls{prototipo}.
    \item \textbf{Calidad bibliográfica}: papers, tesis y libros de
          editoriales serias, con datos completos. Wikipedia y blogs restan.
    \item \textbf{Entrevistas a especialistas} (si las hay): fecha, nombre,
          cargo y aporte; transcripción en el apéndice~\ref{apx:entrevista}.
\end{enumerate}
Errores que la cátedra marca: falta de coherencia interna, bibliografía
superficial, información irrelevante o redundante. Es uno de los
capítulos más extensos (orientativo: 20 a 30 páginas).
\end{guia}

La organización de eventos, tanto empresariales como sociales, ha
evolucionado en los últimos años debido al crecimiento de la propia
industria y a la incorporación de nuevas tecnologías. Asimismo, tanto los
clientes como los invitados han aumentado sus expectativas con el paso
del tiempo, siendo cada vez más exigentes respecto a la calidad de los
servicios y a la experiencia que esperan vivir.

Organizar un evento implica coordinar una gran cantidad de actores, entre
ellos salones, servicios de catering, fotógrafos, decoradores, empresas
de entretenimiento, proveedores de equipamiento y distintos profesionales
que deben trabajar de manera coordinada para garantizar el éxito del
evento. Esta coordinación requiere una planificación constante y un
seguimiento de cada una de las tareas que forman parte de la
organización.

A pesar de estos avances, gran parte de la organización continúa
llevándose a cabo mediante plataformas independientes, como redes
sociales, aplicaciones de mensajería y distintos sistemas de gestión.
Esto provoca que los clientes no tengan toda la información de su evento
centralizada en un único lugar, dificultando el seguimiento de las tareas
y aumentando la posibilidad de cometer errores o perder información
importante durante la planificación.

A partir de esta situación, resulta necesario analizar las tecnologías y
herramientas que actualmente se utilizan para la organización de eventos,
así como las distintas plataformas e investigaciones desarrolladas en
este ámbito, con el fin de identificar oportunidades de mejora y
fundamentar el desarrollo de la propuesta planteada en este trabajo.

\section{Concepto de la industria de los eventos}

La industria de los eventos ha crecido mucho en los últimos años,
generando mayor importancia en el ámbito social, cultural y profesional.
Gracias a este crecimiento, distintos organismos e investigadores
desarrollaron definiciones y estudios con el objetivo de comprender cómo
funciona la industria, sus características y los procesos que intervienen
en la gestión y planificación de los eventos.

Donald Getz, uno de los autores más influyentes en el desarrollo del
Event Management y los Event Studies, sostiene en su obra \emph{Event
Studies: Theory, Research and Policy for Planned Events}
\citep{getz2007} que los eventos
planificados constituyen acontecimientos temporales creados con un
propósito determinado, que poseen un inicio y un final claramente
definidos y requieren un proceso de planificación, organización y
gestión para alcanzar los objetivos propuestos.

Joe Goldblatt, uno de los pioneros en la profesionalización de la gestión
de eventos, define a un evento como un momento único en el tiempo,
celebrado mediante ceremonias o actividades específicas con el propósito
de satisfacer determinadas necesidades \citep{goldblatt2002}.

El Events Industry Council (EIC) \citep{eic} considera que esta industria
comprende el conjunto de actividades relacionadas con la planificación,
organización y realización de congresos, convenciones, reuniones,
exposiciones, ferias y otros eventos, involucrando a una amplia red de
organizaciones y profesionales.

Todas las definiciones coinciden en que la organización de eventos
constituye una actividad compleja que requiere planificación,
coordinación y gestión de múltiples recursos, por lo que resulta
fundamental disponer de herramientas tecnológicas que contribuyan a
centralizar la información y optimizar la \gls{g:gestion-integral-eventos}
de cada evento.

\section{Clasificación de los eventos}

\begin{itemize}
    \item \textbf{Eventos sociales}: casamientos, cumpleaños,
          aniversarios, bautismos, fiestas de quince años y reuniones
          privadas, donde la personalización de los servicios es un
          aspecto fundamental.
    \item \textbf{Eventos corporativos}: congresos, convenciones,
          lanzamientos de productos, reuniones empresariales,
          capacitaciones y actividades de integración de equipos,
          orientados al cumplimiento de objetivos organizacionales.
    \item \textbf{Eventos culturales}: recitales, festivales, obras
          teatrales y exposiciones, con requerimientos logísticos ligados
          a venta de entradas, seguridad e infraestructura técnica.
    \item \textbf{Eventos deportivos}: competencias, torneos y
          campeonatos que requieren planificación de instalaciones,
          cronogramas, acreditaciones y seguridad.
    \item \textbf{Ferias, exposiciones y congresos}: orientados al
          intercambio de conocimientos y la generación de negocios, con
          administración simultánea de expositores, asistentes y
          espacios.
\end{itemize}

\section{Principales actores involucrados}

Un evento no depende de una única persona o empresa, sino de la
participación de diferentes actores que aportan servicios, recursos e
información. Entre ellos se destacan:

\begin{itemize}
    \item \textbf{El cliente}: persona, empresa o institución que
          solicita el evento y comunica sus necesidades (tipo de evento,
          fecha, presupuesto, cantidad de invitados).
    \item \textbf{El organizador o coordinador}: actúa como punto de
          conexión entre el cliente, el salón y los proveedores,
          relevando requerimientos, controlando presupuestos y
          supervisando el cumplimiento de cada tarea.
    \item \textbf{Los salones y espacios}: administran su calendario de
          reservas y comunican las condiciones de contratación,
          capacidad e instalaciones disponibles.
    \item \textbf{Proveedores gastronómicos}: planifican la elaboración
          de alimentos y adaptan los menús a las necesidades del cliente.
    \item \textbf{Proveedores de ambientación y equipamiento}:
          decoración, mobiliario, mantelería y arreglos florales,
          coordinando horarios de montaje.
    \item \textbf{Proveedores técnicos}: sonido, iluminación, pantallas y
          equipos audiovisuales.
    \item \textbf{Proveedores audiovisuales y de entretenimiento}:
          fotografía, filmación, DJs y animadores, que requieren conocer
          el cronograma completo de actividades.
\end{itemize}

\section{Proceso general de organización de un evento}

\begin{enumerate}
    \item \textbf{Definición de requerimientos}: identificación de
          objetivos, tipo de celebración, fecha, presupuesto y cantidad
          de asistentes.
    \item \textbf{Búsqueda y selección de proveedores}: evaluación de
          disponibilidad, experiencia, calidad y costos, habitualmente a
          través de redes sociales, sitios web o plataformas de
          \gls{g:marketplace}.
    \item \textbf{Cotización y contratación}: solicitud de presupuestos,
          negociación de condiciones y formalización de reservas
          mediante contratos y señas.
    \item \textbf{Planificación y coordinación}: definición del
          cronograma, distribución del espacio e intercambio permanente
          de información entre proveedores.
    \item \textbf{Seguimiento y control}: verificación de confirmaciones,
          pagos pendientes y cumplimiento de tareas, gestionando cambios
          de último momento.
    \item \textbf{Cierre y evaluación}: devolución de equipamiento,
          cierre de contratos y encuestas de satisfacción para
          identificar mejoras futuras.
\end{enumerate}

\section{Problemáticas durante la organización de eventos}

\begin{itemize}
    \item \textbf{Descentralización de la información}: los datos se
          dispersan entre correos, WhatsApp, redes sociales y planillas,
          dificultando saber cuál es la versión más actualizada.
    \item \textbf{Dificultades en la búsqueda y contratación de
          proveedores}: cada proveedor utiliza formatos distintos para
          su información, lo que exige consultas individuales y una alta
          inversión de tiempo en comparar alternativas.
    \item \textbf{Problemas de coordinación entre actores}: el uso de
          herramientas distintas por cada participante aumenta las
          posibilidades de errores y demoras.
    \item \textbf{Ausencia de una gestión integral}: las herramientas
          existentes resuelven solo una parte del proceso, sin ofrecer un
          seguimiento completo de la organización.
\end{itemize}

\section{Transformación digital en la industria de los eventos}

La pandemia de COVID-19 impulsó la realización de eventos virtuales e
híbridos, acelerando la adopción de plataformas más completas capaces de
centralizar múltiples actividades. Hoy existen soluciones que integran
gestión de clientes, cronogramas, proveedores, presupuestos y generación
de indicadores. Sin embargo, la mayoría de las empresas del sector
todavía no alcanzó un nivel de madurez digital suficiente para
aprovechar plenamente la \gls{ia}, cuyo uso suele limitarse a asistentes
virtuales o recomendaciones básicas, sin una integración completa en la
planificación de eventos.

\section{Plataformas existentes para la gestión y organización de eventos}\label{sec:plataformas-existentes}

\subsection{Internacionales}

Cvent es una de las plataformas más utilizadas para la gestión integral
de eventos corporativos, congresos y convenciones, con funcionalidades
basadas en \gls{ia} orientadas a eventos empresariales de mediana y gran
escala. Bizzabo se especializa en congresos y eventos corporativos, con
herramientas de personalización de la experiencia de los asistentes.
Whova ofrece gestión de asistentes, agendas y networking para congresos y
reuniones empresariales, aunque no cuenta con herramientas de gestión
integral de proveedores ni de simulación. WeddingWire y The Knot
funcionan como marketplaces especializados en la organización de bodas,
permitiendo localizar proveedores y consultar opiniones, aunque la
gestión posterior del evento continúa realizándose mediante herramientas
externas.

\begin{table}[H]
    \centering
    \caption{Plataformas internacionales relevadas}
    \label{tab:plataformas-internacionales}
    \small
    \begin{tabularx}{\textwidth}{l L c c c c}
        \toprule
        \tb{Plataforma} & \tb{Enfoque principal} & \tb{Marketplace} & \tb{Gestión integral} & \tb{IA} & \tb{Gemelo Digital} \\
        \midrule
        Cvent & Eventos corporativos, congresos y convenciones & No & Sí & Parcial & No \\
        Bizzabo & Congresos y eventos corporativos & No & Sí & Parcial & No \\
        Whova & Congresos y reuniones empresariales & No & Parcial & No & No \\
        WeddingWire & Organización de bodas y contratación de proveedores & Sí & No & No & No \\
        The Knot & Organización de bodas y contratación de proveedores & Sí & No & No & No \\
        \bottomrule
    \end{tabularx}
\end{table}

\subsection{Nacionales}

BIG Fiestas y FEX Eventos son plataformas argentinas orientadas a la
organización de eventos sociales mediante un modelo de
\gls{g:marketplace}, centralizando la búsqueda de proveedores sin
incorporar herramientas de gestión integral. Casamientos.com.ar es una de
las plataformas argentinas más reconocidas para bodas, enfocada
principalmente en la contratación de servicios.

\begin{table}[H]
    \centering
    \caption{Plataformas nacionales relevadas}
    \label{tab:plataformas-nacionales}
    \small
    \begin{tabularx}{\textwidth}{l L c c c c}
        \toprule
        \tb{Plataforma} & \tb{Enfoque principal} & \tb{Marketplace} & \tb{Gestión integral} & \tb{IA} & \tb{Gemelo Digital} \\
        \midrule
        BIG Fiestas & Organización de eventos sociales y contratación de proveedores & Sí & No & No & No \\
        FEX Eventos & Búsqueda y contratación de proveedores & Sí & No & No & No \\
        Casamientos.com.ar & Organización de bodas y contratación de proveedores & Sí & No & No & No \\
        \bottomrule
    \end{tabularx}
\end{table}

\subsection{Brecha identificada}

Ninguna de las plataformas relevadas, nacionales ni internacionales,
combina las tres capas que propone este proyecto: \gls{g:marketplace} de
servicios, gestión integral (reservas, presupuestos, contratos, pagos) y
funcionalidades inteligentes de recomendación y simulación. Las
plataformas de gestión corporativa (Cvent, Bizzabo) resuelven la gestión
pero no operan como marketplace abierto; los marketplaces de bodas
(WeddingWire, The Knot, BIG Fiestas, FEX Eventos, Casamientos.com.ar)
resuelven la búsqueda y contratación pero no ofrecen herramientas de
gestión posterior ni de simulación. Esta brecha es la que justifica el
desarrollo de una plataforma integral en lugar de sumarse a una categoría
ya cubierta por la oferta existente.

\section{Inteligencia artificial aplicada a la organización de eventos}

La \gls{ia} puede definirse como un conjunto de tecnologías capaces de
procesar grandes volúmenes de información para identificar patrones,
generar recomendaciones y asistir en la toma de decisiones. En este
proyecto se busca que el sistema analice la información disponible para
sugerir las opciones más compatibles con las necesidades del usuario,
considerando variables como disponibilidad, capacidad, precio, proximidad
geográfica y valoraciones de otros clientes.

\subsection{¿Qué es un Gemelo Digital?}

El \gls{g:gemelo-digital} (Digital Twin) es una representación virtual de
un objeto, proceso o sistema real que permite analizar su comportamiento
mediante información proveniente del entorno que representa. A diferencia
de un modelo estático o de una simulación tradicional, el Gemelo Digital
se actualiza continuamente con nuevos datos, permitiendo simular
distintos escenarios a medida que cambia la información del sistema real.

\begin{table}[H]
    \centering
    \caption{Modelo estático, simulación tradicional y Gemelo Digital}
    \label{tab:gemelo-digital-comparacion}
    \small
    \begin{tabularx}{\textwidth}{L L L}
        \toprule
        \tb{Modelo digital estático} & \tb{Simulación tradicional} & \tb{Gemelo Digital} \\
        \midrule
        Representación digital de un objeto o proceso. No cambia
        automáticamente ni analiza comportamientos por sí solo. &
        Modelo que ejecuta uno o varios escenarios definidos. Se ejecuta
        para analizar un escenario específico. &
        Representación digital que se actualiza con datos y permite
        simulaciones continuas. Evoluciona junto con el sistema real. \\
        \bottomrule
    \end{tabularx}
\end{table}

Aplicado al proyecto, el Gemelo Digital representa virtualmente cada
evento registrado por el usuario (tipo, fecha, ubicación, invitados,
presupuesto y proveedores). A medida que el usuario modifica esta
información, el modelo actualiza su estado, permitiendo evaluar el
impacto de decisiones como un aumento en la cantidad de invitados sobre
la capacidad del salón, el catering o el presupuesto.

\section{Marco teórico: relación de variables del proyecto}

Siguiendo el enfoque planteado por la cátedra, se identifican a
continuación las dos variables centrales sobre las que se apoya el
proyecto, explicando cómo se relacionan entre sí en lugar de tratarlas
como conceptos aislados.

\subsection{Variable 1: plataformas de marketplace (economía de plataformas)}

Un \gls{g:marketplace} digital constituye lo que la literatura económica
denomina \gls{g:economia-plataformas} (\emph{two-sided market}): una
plataforma que conecta a dos grupos de usuarios distintos ---clientes que
buscan contratar servicios de eventos y salones/proveedores que ofrecen
esos servicios--- cuyo valor para cada lado depende de la cantidad y
calidad de participantes del lado opuesto. \citet{rochettirole2006}
formalizaron este concepto y explicaron por qué este tipo de plataformas
suele requerir estrategias de precios diferenciadas para cada lado del
mercado.

\subsection{Variable 2: sistemas de recomendación basados en inteligencia artificial}

Los \gls{g:sistema-recomendacion} procesan información sobre
preferencias, comportamiento y características de los usuarios para
sugerir automáticamente las opciones más adecuadas. \citet{riveroetal2025},
en una revisión de literatura sobre automatización mediante \gls{ia} en
el ámbito turístico, identificaron que la principal tendencia de
investigación actual es el desarrollo y la personalización de sistemas de
recomendación, confirmando la relevancia de este tipo de herramientas en
industrias de servicios con alto grado de intermediación.

\subsection{Relación entre ambas variables}

En un mercado de dos lados con una gran cantidad de proveedores y
combinaciones posibles, como el de la organización de eventos, el costo
de búsqueda que enfrenta el cliente crece rápidamente a medida que
aumenta la oferta disponible. Un sistema de recomendación basado en
\gls{ia} reduce ese costo de búsqueda, priorizando las combinaciones con
mayor compatibilidad con las características del evento. Esta relación
no es unidireccional: cuantos más usuarios utiliza el marketplace, más
datos genera, y esos datos son el insumo principal para mejorar la
precisión de las recomendaciones, en un ciclo de retroalimentación entre
ambas variables. El Gemelo Digital se apoya sobre esta misma relación,
utilizando tanto los datos del marketplace como las salidas del sistema
de recomendación para simular el impacto de distintas decisiones antes de
que el evento se lleve a cabo.

\section{Marco conceptual}

A continuación se definen conceptos utilizados a lo largo del trabajo que
no fueron desarrollados en profundidad en las secciones anteriores.

\begin{description}
    \item[\gls{g:economia-plataformas}] Modelo de negocio digital en el
          que una plataforma crea valor principalmente al facilitar las
          interacciones entre dos o más grupos de usuarios
          interdependientes \citep{rochettirole2006}.
    \item[\gls{ml}] Subcampo de la \gls{ia} que permite que un sistema
          mejore su desempeño en una tarea a partir del análisis de
          datos, sin ser programado explícitamente para cada caso.
    \item[\gls{g:big-data}] Conjuntos de datos de gran volumen, velocidad
          y variedad, base sobre la que operan los sistemas de
          recomendación y los Gemelos Digitales.
    \item[\gls{g:cloud-computing}] Modelo de provisión de infraestructura,
          almacenamiento y procesamiento a través de Internet, sin
          administrar servidores físicos propios.
    \item[\gls{g:mvp}] Versión mínima de un producto, funcional y
          utilizable por usuarios reales, que permite validar una
          propuesta de valor con la menor inversión de desarrollo
          posible.
    \item[\gls{api} REST] Conjunto de reglas que permite la comunicación
          entre el frontend y el backend de una aplicación mediante
          solicitudes HTTP estandarizadas.
    \item[\gls{stakeholder}] Toda persona, grupo u organización que
          participa del proyecto o se ve afectada por sus resultados.
\end{description}

\section{Antecedentes de investigación: estudios nacionales e internacionales}

Además de las plataformas comerciales relevadas en la
sección~\ref{sec:plataformas-existentes}, se identificaron los siguientes
antecedentes académicos vinculados a las variables centrales del
proyecto:

\begin{itemize}
    \item \citet{licht2023} --- Universidad de Palermo (tesis de
          maestría, Argentina). Plan de negocios para un marketplace que
          conecta productores de frutas y verduras con consumidores en la
          Ciudad de Buenos Aires, con un modelo de ingresos mixto
          (suscripción + comisión del 5\,\%) evaluado mediante \gls{van}
          y \gls{tir}, validando la aplicabilidad de esta metodología a
          proyectos de marketplace en Argentina.
    \item \citet{riveroetal2025} --- Universidad Nacional de La Plata.
          Revisión de literatura sobre automatización mediante \gls{ia}
          en el ámbito turístico, que concluye que la principal línea de
          investigación actual es la personalización de sistemas de
          recomendación.
    \item \citet{chicaizaetal2024} --- Universidad de Sevilla. Revisión
          histórica del concepto de Gemelo Digital, que señala su
          expansión reciente desde la industria manufacturera hacia otros
          campos, respaldando su aplicación a un dominio no industrial
          como la organización de eventos.
    \item \citet{varaschiquito2020} --- RECIMUNDO. Revisión bibliográfica
          sobre la evolución de los Gemelos Digitales, describiendo las
          tecnologías habilitantes (IoT, Big Data, IA y Cloud Computing)
          que coinciden con las consideradas en el marco tecnológico del
          proyecto.
    \item \citet{rochettirole2006} --- The RAND Journal of Economics.
          Artículo de referencia sobre mercados de dos lados, que
          formaliza el marco teórico utilizado en la
          sección~\ref{sec:plataformas-existentes}.
\end{itemize}

\section{Entrevistas a especialistas}

Se realizaron dos entrevistas semiestructuradas a especialistas del
sector (transcripción completa en el apéndice~\ref{apx:entrevista}): una
al dueño de un salón de eventos con más de diez años de experiencia
(22/07/2026), y otra a la gerente de planificación y ejecución de una
empresa de decoración y ambientación con más de diez años de experiencia
(23/07/2026). Ambas coinciden en identificar la coordinación entre
proveedores y la actualización de la información de pagos como los
puntos más difíciles de sostener con las herramientas actuales
(WhatsApp, Excel, agenda en papel), y se muestran receptivas a
funcionalidades de calendario compartido, registro de pagos y consulta de
información según el rol de cada usuario. La utilidad de las
recomendaciones automáticas y de la simulación previa del evento generó
opiniones más dispares: mientras la entrevista de decoración las valora
positivamente, la de salón se muestra cauta con las recomendaciones y ve
la simulación como una herramienta pensada sobre todo para el cliente.

Queda pendiente la entrevista a un responsable de catering, que falta
para completar la Fase~2 de la metodología.

\section{Encuesta a posibles usuarios}

Para triangular con las entrevistas anteriores, se relevó también a
potenciales clientes de la plataforma mediante una encuesta autoadministrada
(cuestionario completo y metodología en el apéndice~\ref{apx:encuesta}),
abierta entre el 22/09/2026 y el 29/09/2026, con 45 respuestas.

\subsection{Cómo organizan hoy su evento}

El 91,1\,\% de quienes respondieron ya había organizado un evento
(casamiento, cumpleaños, fiesta de 15 u otro), la mayoría de entre 50 y
100 invitados (42,2\,\%). Los medios más usados para buscar y contratar
proveedores fueron Instagram o Facebook y la recomendación de conocidos
(55,6\,\% cada uno), seguidos por WhatsApp (48,9\,\%) y Google
(31,1\,\%); solo un 2,2\,\% usó alguna app o sitio de tipo marketplace
ya existente (BIG Fiestas, Casamientos.com.ar), lo que confirma que
estas plataformas todavía no son el canal principal de búsqueda relevado
en la sección~\ref{sec:plataformas-existentes}.

\begin{table}[H]
    \centering
    \caption{Esfuerzo de búsqueda y contratación reportado}
    \label{tab:encuesta-esfuerzo}
    \small
    \begin{tabularx}{\textwidth}{L L}
        \toprule
        \tb{Proveedores contactados} & \tb{Tiempo dedicado a buscar y comparar} \\
        \midrule
        1 a 3: 40,0\,\% & Menos de 1 semana: 24,4\,\% \\
        4 a 6: 57,8\,\% & Entre 1 y 4 semanas: 46,7\,\% \\
        7 o más: 2,2\,\% & Más de 1 mes: 28,9\,\% \\
        \bottomrule
    \end{tabularx}
\end{table}

Lo más difícil de organizar fue, en orden, comparar precios (64,4\,\%),
encontrar opciones disponibles (37,8\,\%), coordinar horarios entre
proveedores (22,2\,\%) y llevar el control de pagos y presupuesto
(20,0\,\%). En una escala de 1 a 5, la complejidad percibida promedió
3,04, con un 35,6\,\% calificándola 4 o 5.

\subsection{Interés en la propuesta}

El 82,2\,\% usaría ``seguro'' una única plataforma para buscar, comparar
y contratar todos los servicios de su evento, y un 8,9\,\% adicional
``probablemente'' la usaría (91,1\,\% en conjunto). El interés en las
dos funcionalidades inteligentes es todavía mayor: 93,3\,\% valoró útil
que la plataforma recomiende proveedores mediante \gls{ia} según
presupuesto, ubicación y cantidad de invitados, y 97,8\,\% valoró útil
poder simular el evento antes de confirmarlo.

\subsection{Conclusión de la encuesta}

Los resultados respaldan tanto el diagnóstico del
capítulo~\ref{cap:introduccion} (la búsqueda y comparación de
proveedores consume semanas y varios canales distintos, sin que los
marketplaces existentes hayan desplazado a WhatsApp y las redes sociales)
como el interés en las tres capas de la propuesta: marketplace,
recomendación y Gemelo Digital. El interés declarado en la recomendación
y la simulación superó incluso al interés en la plataforma unificada en
sí, lo que sugiere que el componente inteligente no es un agregado
secundario sino parte de lo que motivaría la adopción. Al tratarse de una
muestra por conveniencia y no probabilística (apéndice~\ref{apx:encuesta}),
estos porcentajes se toman como una validación exploratoria de la
dirección del proyecto, no como una medición representativa de toda la
población.
